\clearpage
\section*{Introducere}
\addcontentsline{toc}{section}{Introducere}

Bateriile reprezintă o sursă vitală de energie în numeroase aplicații, de la telefoane mobile și laptopuri, până la mașini electrice și sisteme de stocare a energiei. Cu toate acestea, durata de viață și performanța bateriilor pot fi influențate de diferite factori. De aceea, dezvoltarea și implementarea unui sistem de gestionare a bateriilor eficient devine din ce în ce mai importantă. Un sistem de gestionare a bateriilor are rolul de a monitoriza și controla starea bateriilor, optimizând astfel durata lor de viață, eficiența și performanța. În această lucrare, va fi demonstrată importanța unui astfel de sistem și modul în care acesta poate aduce beneficii semnificative în domeniul energiei și al mobilității electrice.

DIYBMS (Do-It-Yourself Battery Management System) \cite{bms-intro} este un proiect open-source care oferă o soluție de monitorizare și gestionare a bateriilor cu ioni de litiu (Li-Ion), dezvoltat cu scopul de a fi integrat într-un sistem de stocare a energiei regenerabile. Proiectul a fost inițial dezvoltat de către Stuart Pittaway, un inginer din Marea Britanie, și a fost lansat în 2017. 

DIYBMS oferă o serie de posibilități, cum ar fi monitorizarea tensiunii și a temperaturii bateriilor, protecție împotriva supraîncărcării și descărcării excesive, precum și posibilitatea de a interconecta mai multe module DIYBMS pentru a gestiona un număr mare de baterii. Cu toate acestea, există o limitare în ceea ce privește tensiunea bateriilor pe care le poate suporta DIYBMS, aceasta nu trebuie să depășească tensiunea de 5V.

Pentru a depăși această limitare, sistemul DIYBMS a fost adaptat pentru a permite gestionarea și monitorizarea bateriilor de tip Nichel-Metal-Hidrid (NiMH). Aceste baterii au o tensiune nominală de 1,2 volți pe celulă și sunt alcătuite din 6 celule, rezultând o tensiune totală de 7,2 volți. Modificarea adusă DIYBMS a oferit utilizatorilor posibilitatea de a gestiona și monitoriza bateriile NiMH, evitând astfel necesitatea achiziționării unui sistem de gestionare a bateriilor mai costisitor.

În acest raport, vor fi expus procesul de adaptare al proiectului DIYBMS pentru a răspunde cerințelor de gestionare și monitorizare a bateriilor de tip NiMH. Vor fi ilustrate modificările aduse sistemului pentru a suporta aceste baterii și va fi documentat modul de utilizare al sistemului realizat. Raportul data poate fi folosit și ca instrucțiuni pentru construirea unui astfel de sistem, fiind că acesta nu poate fi achiziționat de la vreun producător.
